\chapter{Implementation}
\label{ch:implementation-chapter}
\newpage

This chapter details the concrete technical implementation of the NewsPulse microservices architecture. It transitions the structural diagrams and logical designs established in previous phases into functional software systems, detailing the programmatic flows, exact configuration boundaries, and specific third-party integration pipelines utilized to build the system. By rejecting expensive and high-latency cloud architectures in favor of highly optimized local components, the implementation prioritizes resource-constrained efficiency. It integrates CPU-bound machine learning pipelines with deterministic algorithmic fallbacks to deliver low-latency news intelligence on standard consumer hardware.

\section{Important Flow Control and Algorithmic Logic}
The functional stability of the platform is governed by explicit algorithmic flow control loops, engineered to enforce performance limits and prevent computational resource exhaustion.

\subsection{Asynchronous Scheduled Ingestion Loop}
The system's primary ingestion cycle is operated by the automated background aggregation daemon within the \texttt{article\_fetcher} container. Rather than running a blocking linear script, this daemon leverages an asynchronous lifespan environment within the FastAPI framework. The startup sequence invokes a dedicated non-blocking polling routine managed via an asynchronous scheduling loop executing every 15 minutes. 

To bypass anti-bot limits, incomplete feeds, and Transport Layer Security (TLS) configuration restrictions, the ingestion pipeline implements a rigorous three-tier degradation crawler:
\begin{itemize}
    \item \textbf{Tier 1 (Spoofed HTTPX Requests):} The daemon initiates an asynchronous request to the configured NewsAPI, NewsData.io, or GNews endpoints using randomized user-agent headers to simulate native web browser headers.
    \item \textbf{Tier 2 (Trafilatura Native Extraction):} If the upstream API returns an incomplete payload or truncated text body, the system shifts to a secondary scraper that spawns a \texttt{trafilatura} extraction instance. This tool parses the raw Hyper-Text Markup Language (HTML) dynamically to isolate the primary news text from boilerplate scripts, sidebars, and navigation elements.
    \item \textbf{Tier 3 (Bypassed TLS Sockets):} If strict local security certificates block standard secure handshakes, the pipeline degrades to a direct socket connection utilizing the \texttt{urllib3} pool manager. This manager is explicitly configured with unverified TLS settings (\texttt{cert\_reqs='CERT\_NONE'}) to force payload recovery.
\end{itemize}
To protect the system's storage and downstream analytics from processing noise, a post-scraping filter discards any extracted texts containing fewer than 100 characters. Furthermore, computationally heavy HTML parsing routines are offloaded to a designated local \texttt{ThreadPoolExecutor} to prevent blocking the primary asynchronous event loop.

\subsection{Veracity Sentinel Consensus Algorithm}
To evaluate news credibility without high-overhead stochastic models, the veracity engine executes a multi-signal consensus formula inside the \texttt{fakenews\_detection} service. The consensus score incorporates four distinct signals:
\begin{enumerate}
    \item \textbf{Supervised Inference (45\%):} Sequence classification from the CPU-optimized \texttt{dhruvpal/fake-news-bert} model.
    \item \textbf{StylometricSensationalism Heuristic (25\%):} Capitalization frequency ratios, exclamation mark density, and regex-based sensational clickbait patterns.
    \item \textbf{Cognitive Complexity Heuristic (20\%):} Lexical diversity (unique token ratios) and factual vs. emotional vocabulary density.
    \item \textbf{Baseline Anchor (10\%):} A neutral mathematical anchor set at 0.5.
\end{enumerate}
At runtime, the service calculates the weighted sum of these four factors. The resulting raw score is then adjusted by a domain-specific credibility multiplier:
\begin{equation}
    P_{\text{adjusted}} = P_{\text{raw}} \times M_{\text{domain}}
\end{equation}
where the multiplier $M_{\text{domain}}$ is set to $0.8$ for verified, trusted publishers and $1.3$ for unverified sources. 

If the title matches verified facts checked via live APIs, the fake news probability is immediately scaled down to 25\% of the BERT classifier score. If the style analyzer detects a highly objective document (stylometric sensationalism $< 0.1$ and emotional density $< 0.3$), a hard cap overrides the final probability to a maximum of $0.49$. This ceiling suppresses false positive machine learning flags. The engine also extracts the sentence with the highest sensational density to serve as the explainable highlight phrase in the client interface.

\subsection{Sliding Window Bias Classification and Topic Gating}
The political bias container runs a local instance of the \texttt{valurank/distilroberta-bias} model. Because standard RoBERTa architectures are limited to a maximum sequence length of 512 tokens, long-form articles are evaluated using an overlapping sliding window algorithm. The raw text is divided into token arrays using a sequence window of 512 tokens, a stride of 448 tokens, and an overlap of 64 tokens.

The final polarization vector is calculated by computing the weighted average of the chunk probabilities. The system applies a 1.5x weight multiplier to the introduction and conclusion chunks, as these sections historically contain the strongest framing signals. To prevent false positives in apolitical texts, the classification is governed by an apolitical keyword penalty gate. If an article contains technical, scientific, or athletic terms without at least two political keywords, a penalty value between $0.3$ and $0.7$ overrides the classification, forcing the final label to "Center" with high confidence.

\section{Components, Libraries, Web Services and Stubs}
The NewsPulse system is built using a highly decoupled Python and React software stack, where each package is strictly isolated within its respective container.

\subsection{Core Frameworks and Persistence}
\begin{itemize}
    \item \textbf{FastAPI:} Serves as the high-throughput asynchronous web framework that powers all backend services and the primary API Gateway.
    \item \textbf{SQLAlchemy ORM:} Manages relational databases within the authentication and analytics logging microservices, utilizing connection pools optimized for SQLite.
    \item \textbf{Passlib bcrypt:} Implements cryptographically secure, salted bcrypt hashing within the user management service to protect user credentials.
\end{itemize}

\subsection{Machine Learning and Transformers}
The localized natural language processing pipelines leverage the HuggingFace \texttt{transformers} library configured strictly for CPU execution. Specific models include:
\begin{itemize}
    \item \textbf{Abstractive Summarizer:} The \texttt{sshleifer/distilbart-cnn-12-6} model, chosen for its low memory footprint and high compression ratio on CPU hardware.
    \item \textbf{Credibility Classifier:} The \texttt{dhruvpal/fake-news-bert} model, utilized to output sequence classification probabilities.
    \item \textbf{Bias Detector:} The \texttt{valurank/distilroberta-bias} model, deployed to identify left, center, and right political framing.
    \item \textbf{Socratic Debater:} The \texttt{MBZUAI/LaMini-Flan-T5-248M} model, quantized to INT8 to generate critical counter-arguments on standard processors.
\end{itemize}

\subsection{Cryptographic Security and API Stubs}
To secure data in transit, the system implements a hybrid asymmetric-symmetric cryptosystem. The Python standard \texttt{cryptography} package on the backend and the \texttt{node-forge} library in the React frontend handle the encryption routines. 
\begin{itemize}
    \item \textbf{Asymmetric Layer:} Encapsulates transient keys using the RSA-2048 standard with Optimal Asymmetric Encryption Padding (OAEP) utilizing the SHA-256 hash function.
    \item \textbf{Symmetric Layer:} Encrypts the data payload using the 32-byte AES standard in Galois/Counter Mode (AES-256-GCM) with an initialization vector (IV) and authentication tag.
    \item \textbf{OAuth and Stubs:} The React SPA manages third-party authentication using the \texttt{@react-oauth/google} package. The ingestion layer connects to external news APIs (NewsAPI, NewsData.io, and GNews) as stubs to feed the background aggregation loop.
\end{itemize}

\section{Deployment Environment}
The system is deployed using a decoupled containerized environment configured via Docker Compose to maintain exact parity across development, testing, and production states.

\begin{figure}[h!]
    \centering
    \includegraphics[width=1\linewidth]{BS_FYP_Figs/deployement_diagram.png}
    \caption{Docker Compose Deployment Topology and Network Bridging}
    \label{fig:deployment_topology}
\end{figure}

As shown in Figure~\ref{fig:deployment_topology}, the production backend is hosted on localized consumer off-the-shelf (COTS) server hardware. The network topology relies on two core configurations:
\begin{itemize}
    \item \textbf{Docker Compose Private Bridge Network:} The backend stack comprises 11 decoupled containers isolated within a private bridge network. The only port exposed to the host machine is the API Gateway (Port 8000). A \texttt{Redis 7.4-alpine} container is deployed within this private network, acting as the primary hot-caching database to eliminate redundant inferences.
    \item \textbf{Secure Transit Tunneling:} An outbound HTTPS tunnel is established on the host using \texttt{ngrok} or Cloudflare Tunnels. This tunnel binds the gateway port (8000) to a public HTTPS URL. This allows the client application to communicate securely with the gateway without requiring public IP addresses or exposing administrative ports to the public internet. The React application is hosted on Vercel's global static edge network.
\end{itemize}

\section{Tools and Techniques}
To achieve the No-LLM "Socratic Debater Engine" and ensure system stability under heavy workloads, the platform implements advanced natural language processing optimization techniques.

\subsection{PyTorch Dynamic INT8 Quantization}
The sequence-to-sequence T5 model is highly memory-intensive when executing on standard CPUs. To optimize inference speeds, the system applies PyTorch dynamic quantization (\texttt{torch.quantization.quantize\_dynamic}) to the model's linear layers during container initialization. This process converts 32-bit floating-point weights into highly compressed 8-bit integers:
\begin{equation}
    W_{\text{quantized}} = \text{round}\left( \frac{W_{\text{float32}}}{\text{Scale}} \right) + \text{ZeroPoint}
\end{equation}
This quantization reduces the model's memory footprint by approximately 75\% and improves inference speeds by over 3x on CPU architectures without degrading the logical quality of the output.

\subsection{Heuristic Debater Post-Processing and Rule-Based Fallbacks}
To clean output from smaller language models, the debater container processes the generated text through a multi-stage filtering pipeline:
\begin{itemize}
    \item \textbf{Length and Fragment Filters:} Sentences containing fewer than 20 characters or fewer than 5 words are discarded as fragments.
    \item \textbf{Run-On and Repetition Filters:} Sentences exceeding 40 words are discarded. The system also tracks token frequencies; if a substantive word (length $> 3$) is repeated more than 3 times within a single sentence, it is discarded to prevent model looping.
    \item \textbf{Dangling Preposition Filter:} The system drops any sentence ending with dangling conjunctions or prepositions (e.g., "and", "but", "with", "for", "of").
    \item \textbf{Negative Density Filter:} Sentences containing two or more negations (e.g., "not") are discarded to prevent convoluted or self-contradicting arguments.
\end{itemize}

If a pipeline failure or timeout occurs (using a hard 90-second timeout), the debater executes a rule-based fallback engine. This fallback uses regular expressions to extract proper nouns from the article text and identifies the domain. The extracted proper noun is then dynamically injected into domain-specific templates to generate a structured Socratic response, completely eliminating the risk of AI hallucination.

\subsection{TF-IDF Extractive Summarization Fallback}
If the summarizer service experiences localized memory pressure or pipeline failure, it instantly transitions to a deterministic Term Frequency-Inverse Document Frequency (TF-IDF) extraction algorithm. This fallback tokenizes the article into sentences, filters out common stop-words, and calculates a TF-IDF weight vector for all remaining tokens. Sentences are scored based on the sum of their word weights normalized by the square root of the sentence length:
\begin{equation}
    \text{Score}(S) = \frac{\sum_{w \in S} \text{TF-IDF}(w)}{\sqrt{\text{Length}(S)}}
\end{equation}
The service then extracts and returns the top five highest-scoring sentences to maintain summarization availability.

\section{Best Practices and Coding Standards}
The architecture enforces strict enterprise software engineering standards to ensure high reliability, comprehensive system visibility, and localized security.

\subsection{Asynchronous Multiprocess Telemetry Tracking}
To collect operational telemetry in a multi-worker production environment, all FastAPI services expose standard \texttt{/metrics} endpoints. The services use the \texttt{prometheus\_client} library configured with a shared multiprocess directory. When Uvicorn or Gunicorn spawns multiple parallel worker processes, standard in-memory metrics gatherers fail due to inter-process isolation. 

The system implements the \texttt{MultiProcessCollector} to write worker metrics to a shared, high-speed temporary file directory. The Prometheus scraper polls this endpoint, allowing the system to aggregate request latencies and cache hit ratios without thread lock issues.

\subsection{Shared-Nothing Database Isolation}
To prevent database lock contention and ensure absolute system decoupling, the architecture rejects centralized monolithic databases. Each microservice operates a completely isolated SQLite database mounted to its respective container volume. Cross-service operations are coordinated via asynchronous REST communication and verified using JWT signatures, preventing cascading database failures and locking issues.

\subsection{Google Identity Token Validation Skew}
Containerized environments running on localized Virtual Machines (such as Docker Desktop on Windows host machines) can suffer from clock desynchronization due to host-system power saving states. When validating third-party Google OAuth identity tokens, strict cryptographic systems will immediately reject valid tokens if the system clock drifts even slightly. To handle this, the gateway's token validation handler implements a clock skew allowance of 10 seconds. This skew window handles virtual machine desynchronization while maintaining strong security standards.

\section{Version Control}
The project source code is managed using Git for distributed version control. Development streams are separated into distinct branches (such as authentication, analytics, and crawler pipelines), ensuring developer isolation. Functional code features are merged into the stable \texttt{main} branch only after successfully passing container integration tests and security parameter validations.
